\documentclass[a4paper]{article}

\usepackage{graphicx}
\usepackage{amsmath,amssymb}
\usepackage{minted}
\usepackage{multirow}
\usepackage{float}
\usepackage{algorithm}
\usepackage{algpseudocode}
\usepackage{todonotes}
\usepackage{forest}
\usepackage{xstring}
\usepackage{xcolor}
\usepackage[most]{tcolorbox}
\usepackage{xparse}
\usepackage{natbib}

\title{Large-document startup and Structure responsiveness}
\date{2026-09-03}
\author{Mohammad Rahmani}

\begin{document}
\maketitle

\section{Bug}
A multi-megabyte LaTeX file restored as the active tab could leave the Mind
Mirror window black and unresponsive until it was force-quit. The same file
could also make an explicit open operation expensive even though Visual mode
was already windowed.

\section{Cause}
Three independent full-document costs occurred during Source opening:
\begin{itemize}
    \item LaTeX syntax highlighting evaluated every regular-expression rule on every prose block.
    \item Structure parsing repeatedly counted newlines from the beginning of the source for each structural item, producing quadratic growth.
    \item The initial complete Structure parse ran on the GUI thread.
\end{itemize}

\section{Fix}
The LaTeX highlighter now uses cheap literal triggers and skips regular
expressions for large plain-prose blocks. The Structure parser builds one
line-start index and reuses it for all headings and images. Large LaTeX
Structure parsing is executed by \texttt{LatexStructureParseWorker} outside
the GUI thread.

The complete Source string remains canonical. The Structure worker receives
the complete Source and returns the complete hierarchy; only the timing of the
parse changed. Visual mode continues to use its separate bounded projection.

The stuck session was also recovered by removing the problematic document
from the restored open/active tab state without modifying the document itself.

\end{document}
